% =============================================================================
% ISI Country Brief — Latvia
% External Supplier Concentration Profile
%
% ISI Paper Series — Country Brief
% International Sovereignty Institute — Research Unit
% March 2026
% =============================================================================
\documentclass[12pt,a4paper]{article}

% --- ISI Paper Series shared template ----------------------------------------
\usepackage{isi-series}

% --- PDF metadata ------------------------------------------------------------
\hypersetup{
  pdftitle={External Supplier Concentration Profile: Latvia — ISI EU-27 2024},
  pdfauthor={International Sovereignty Institute -- Research Unit},
  pdfcreator={International Sovereignty Institute},
  pdfsubject={ISI Paper Series -- Country Brief LV, EU-27, Vintage 2024},
  pdfkeywords={sovereignty index, EU-27, supplier concentration,
    latvia, country brief},
}

% ==============================================================================
\begin{document}
\hypersetup{pageanchor=false}

% ---------- Title Page with Metadata ------------------------------------------
\begin{titlepage}
\centering
\vspace*{2cm}
{\Large\textsc{ISI Paper Series --- Country Brief}}\\[1.2cm]
{\LARGE\bfseries External Supplier Concentration Profile:\\[0.3em]
Latvia}\\[1.5cm]
{\large International Sovereignty Index (ISI)\\[0.3em]
EU-27 --- Vintage 2024 --- Methodology v1.0}\\[1.2cm]
{\large International Sovereignty Institute --- Research Unit}\\[0.8cm]
{\large March 2026}

\vspace{0.8cm}
\noindent\rule{\textwidth}{0.4pt}
\vspace{0.4cm}

% --- Research Metadata --------------------------------------------------------
\begin{flushleft}
\noindent\textbf{Data basis:} ISI Paper~2 --- EU-27 Empirical Results 2024 (Methodology~v1.0, Vintage~2024).\\[4pt]
\noindent\textbf{Methodology:} ISI Paper~1 --- Methodological Foundations, Version~1.0.

\medskip

\noindent\textbf{JEL Codes:} F02, F15, F52, O30, Q43\\[4pt]
\noindent\textbf{Keywords:} sovereignty index; EU-27; supplier concentration; composite indicator; variance decomposition; defence dependence

\medskip

\noindent\textbf{Related publications in the ISI Paper Series:}
\begin{itemize}[nosep,leftmargin=1.5em]
\item ISI Paper~1 --- Technical Specification (v1.0) --- doi:\\
  \url{https://doi.org/10.5281/zenodo.18764227}
\item ISI Paper~2 --- EU-27 Empirical Results 2024 (v1.0) --- doi:\\
  \url{https://doi.org/10.5281/zenodo.18764170}
\end{itemize}
\end{flushleft}
\vfill
\end{titlepage}

\hypersetup{pageanchor=true}
\pagenumbering{arabic}
\pagestyle{fancy}
\fancyhead[L]{\small\sffamily\textit{ISI Country Brief}}
\fancyhead[R]{\small\sffamily\textit{Latvia}}
\renewcommand{\headrulewidth}{0.4pt}

% ---------- Body --------------------------------------------------------------

\section{Executive Summary}
\label{sec:summary}

Latvia ranks~25th in the EU-27 with a composite \ISI score of~0.247, classified as \emph{mildly concentrated}.  The dominant axes are logistics~(0.368) and energy~(0.365).  The structural profile is characterised as dual-axis concentration (logistics and energy).  Under Dirichlet weight perturbation (10\,000~Monte Carlo draws), Latvia's mean rank is~25.24 with a standard deviation of~0.45; the 5th--95th percentile band spans ranks~25 to~26.


\section{Country Position in the EU-27}
\label{sec:position}

Latvia occupies rank~25 of~27 EU member states (composite score~0.247).  The EU-27 mean composite score is~0.344 with a standard deviation of~0.070.  Latvia's score lies below the EU-27 mean, indicating below-average external supplier concentration.  In the ranking, Latvia is situated between Spain (rank~24; 0.261) and Slovakia (rank~26; 0.236).


\section{Axis Profile}
\label{sec:axis-profile}

\Cref{tab:axis-profile} presents Latvia's scores on all six \ISI axes.

\begin{table}[htbp]
\centering\small
\caption{Latvia's \ISI axis profile.}
\label{tab:axis-profile}
\begin{tabular}{@{}lr@{}}
\toprule
\textbf{Axis} & \textbf{Score} \\
\midrule
    Finance           & 0.123 \\
    Energy            & 0.365 \\
    Technology        & 0.133 \\
    Defence           & 0.341 \\
    Critical Inputs   & 0.151 \\
    Logistics         & 0.368 \\
\bottomrule
\end{tabular}
\end{table}

The highest axis score is logistics~(0.368), which lies below the EU-27 axis mean of~0.518.  The second-highest axis is energy~(0.365), near the EU-27 axis mean of~0.365.  The lowest axis score is finance~(0.123), below the EU-27 mean of~0.148.


\section{Structural Drivers}
\label{sec:drivers}

\begin{sloppypar}
Latvia's composite score is driven by two axes in combination: logistics~(0.368) and energy~(0.365).  The gap between these two axes is~0.003, indicating a dual-axis concentration profile.  Neither axis alone dominates the composite; rather, the combination of elevated scores on both dimensions determines Latvia's position in the ranking.
\end{sloppypar}


\section{Comparative Context}
\label{sec:comparative}

Across the EU-27, the covariance-based variance decomposition shows that defence (35.1\,\%) and logistics (34.3\,\%) together account for 69.4\,\% of total cross-sectional composite variance.  Latvia's defence score~(0.341) lies below the EU-27 defence mean~(0.573), and its logistics score~(0.368) lies below the EU-27 logistics mean~(0.518).  Under Dirichlet weight perturbation, Latvia's rank standard deviation is~0.45, indicating high rank stability.



Latvia is one of three countries classified as mildly concentrated, with a composite of~0.247 that lies 0.097 below the EU-27 mean.  The axis profile is notably flat: all six scores fall within a 0.24-point range (0.123--0.365), and none exceeds its respective EU-27 mean.  This uniformly low profile distinguishes Latvia from countries at similar composite levels whose profiles are shaped by compensating extremes.

The rank standard deviation of~0.45 and the one-position percentile band~(25--26) make Latvia's rank among the most stable in the entire EU-27.  The explanation is structural: with no axis deviating substantially from the EU-27 mean in either direction, the composite is insensitive to axis reweighting.  Latvia's rank is effectively anchored at the bottom of the distribution by the consistent absence of elevated concentration on any dimension.

The gap to Spain~(rank~24) is~0.014, while the distance to the two co-ranked countries at rank~26---Slovakia and France---is~0.011.  Latvia thus occupies a clearly delineated position, separated from the mid-ranking cluster above but close to the two lowest-ranked member states below.


\section{Interpretation Boundary}
\label{sec:interpretation}

The \ISI measures concentration of external suppliers, not sovereignty in a normative or political sense.  High concentration may reflect geography, economic specialisation, or alliance structures.  A high composite score does not imply policy failure; a low score does not imply policy success.  The index provides an empirical measurement basis --- what policymakers derive from it is a separate question.


% ---------- References --------------------------------------------------------
\clearpage
\vspace*{1.5cm}

\begingroup\emergencystretch=2em\hbadness=10000
\section*{References}

\begin{enumerate}[label={[\arabic*]},leftmargin=2em,itemsep=6pt]
\item International Sovereignty Institute.
  \textit{International Sovereignty Index (ISI) technical specification, version~1.0}.
  Zenodo, 2026.\\ISI Paper Series, Paper~1.
  doi:\,\href{https://doi.org/10.5281/zenodo.18764227}{10.5281/zenodo.18764227}.

\item International Sovereignty Institute.
  \textit{International Sovereignty Index (ISI): EU-27 empirical results 2024 (v1.0)}.
  Zenodo, 2026.\\ISI Paper Series, Paper~2.
  doi:\,\href{https://doi.org/10.5281/zenodo.18764170}{10.5281/zenodo.18764170}.
\end{enumerate}
\endgroup

\end{document}
